\subsection{Accertamento della Qualità}
\label{subsec:accertamento_qualità}

\subsubsection{Scopo}
Il processo di accertamento della qualità ha lo scopo di garantire che i prodotti sviluppati e i processi adottati durante il \glossario{ciclo di vita} del progetto soddisfino i requisiti stabiliti e rispettino le aspettative previste.
Questo processo rappresenta un elemento centrale nella gestione del progetto, in quanto assicura il monitoraggio continuo e la verifica della conformità alle specifiche tecniche, agli obiettivi di progetto e alle normative applicabili.
Il processo si basa su principi fondamentali quali tracciabilità, trasparenza e miglioramento continuo.
Utilizza inoltre \glossario{metriche} cioè strumenti oggettivi che permettono di valutare sia l'efficacia dei processi, valutando l'adesione alle \textit{best practices} dell'ingegneria del software, sia la qualità dei prodotti realizzati.
Attraverso una pianificazione accurata e un monitoraggio costante, il processo di accertamento della qualità contribuisce a ridurre i rischi, migliorare le performance e a soddisfare le aspettative degli \glossario{stakeholder}.

\subsubsection{Garanzia della qualità}
Per garantire un efficace accertamento della qualità, il gruppo adotta il \glossario{ciclo di Deming}, noto anche come PDCA (Plan-Do-Check-Act). 
Questo metodo aiuta a migliorare continuamente processi e prodotti, garantendo il rispetto degli standard richiesti.
Il \glossario{ciclo di Deming} è composto da 4 fasi fondamentali:
\begin{itemize}
    \item \textbf{Plan}: in questa fase si definiscono gli obiettivi di qualità e le strategie per raggiungerli. 
    Vengono quindi stabilite le azioni da intraprendere, le risorse e le tempistiche necessarie per la loro implementazione. 
    Inoltre vengono definite le \glossario{metriche} da raccogliere durante la fase successiva che permetteranno poi di determinare la qualità del processo e/o del prodotto.

    \item \textbf{Do}: in questa fase, le attività pianificate vengono eseguite seguendo la specifica fornita dal piano.
    Durante l'esecuzione delle attività vengono raccolti i dati necessari alla fase successiva.
    

    \item \textbf{Check}: questa fase si verificano i dati raccolti durante la fase "Do" per determinare se i prodotti e i processi sono conformi con gli obiettivi di qualità definiti.
    In questa fase vengono evidenziati i successi e gli aspetti ancora da ottimizzare.
    A supporto dell'analisi vengono utilizzati strumenti grafici che permettono di visualizzare l'andamento dei dati e identificare pattern o anomalie.

    \item \textbf{Act}: Durante questa fase i dati ottenuti nelle fasi "Do" e "Check" vengono analizzati per individuare eventuali criticità.
    Queste possono includere problemi, non conformità, inefficienze, opportunità di miglioramento o qualsiasi altro aspetto che abbia portato a risultati inferiori alle aspettative.
    
\end{itemize}

\subsubsection{Standard di riferimento per la qualità di prodotto}
\label{subsubsec:standard_normativo}
Per la valutazione della qualità del software con conseguente stesura delle \glossario{metriche} viene seguita la normativa \glossario{ISO/IEC 9126}, secondo lo standard le qualità sono suddivise nelle seguenti categorie:
\begin{itemize}
    \item \textbf{qualità esterne}: misurano i comportamenti del software durante la sua esecuzione;
    \item \textbf{qualità interne}: si applicano al software non eseguibile, permettono di individuare eventuali problemi che potrebbero influire sulla qualità finale del prodotto prima che sia realizzato il software eseguibile;
    \item \textbf{qualità in uso}: rappresentano il punto di vista dell'utente sul software, consentono di stabilire i seguenti obiettivi:
    \begin{itemize}
        \item \textit{Efficacia}: capacità del software di permettere agli utenti di raggiungere gli obiettivi specificati con accuratezza e completezza;
        \item \textit{Produttività}: capacità del software di permettere agli utenti di spendere una quantità opportuna di risorse in relazione all'efficacia ottenuta;
        \item \textit{Soddisfazione}: capacità del software di soddisfare gli utenti;
        \item \textit{Sicurezza}: capacità del software di rimanere entro livelli accettabili di rischi di danni a persone e apparecchiature.
    \end{itemize}
\end{itemize}
Lo standard normativo stabilisce un modello definendo un set di caratteristiche che consentono di misurare e valutare diversi aspetti della qualità del prodotto software:
\begin{itemize}
    \item \textbf{Funzionalità}: capacità del software di fornire le funzioni adatte a soddisfare le esigenze stabilite;
    \item \textbf{Affidabilità}: capacità del software di mantenere un livello di prestazioni specifico quando usato in date condizioni per un dato periodo di tempo;
    \item \textbf{Efficienza}: capacità del software di fornire adeguate prestazioni relativamente alla quantità di risorse utilizzate;
    \item \textbf{Usabilità}: capacità del software di essere propriamente compreso e utilizzato dall'utente;
    \item \textbf{Manutenibilità}: capacità del software di essere modificato per introdurre migliorie o adattamenti;
    \item \textbf{Portabilità}: capacità del software di essere trasportato da un ambiente di lavoro a un altro.
\end{itemize}

\subsubsection{Notazione Metriche di Qualità}
Ogni metrica è identificata in modo univoco seguendo la notazione: 
\begin{lstlisting}
        M.<Tipo>.<Abbreviazione Nome>
\end{lstlisting}
Dove:
\begin{itemize}
    \item \textbf{M}: indica "metrica"
    \item \textbf{Tipo}: sarà PC per indicare che la metrica misura la qualità di un processo o PR per indicare che la metrica misura la qualità di un prodotto
    \item \textbf{Abbreviazione Nome}: viene inserito l'acronimo basato sul nome completo della metrica
\end{itemize}
\subsubsection{Didascalia Metriche di Qualità}
Le \glossario{metriche} sono descritte tramite i seguenti campi:
\begin{itemize}
    \item \textbf{Notazione}: seguendo le indicazioni sopra elencate;
    \item \textbf{Nome}: nome completo della metrica;
    \item \textbf{Descrizione}: descrizione della metrica;
    \item \textbf{Caratteristiche}: presente unicamente nelle \glossario{metriche} di prodotto, indica a quale caratteristica, tra quelle indicate nella \hyperref[subsubsec:standard_normativo]{sezione relativa allo standard di riferimento}, si riferisce la metrica;
    \item \textbf{Formula di misurazione}: formula per la misurazione del valore della metrica.
\end{itemize}
Verranno poi indicati all'interno del documento "Piano di Qualifica" i valori tollerabili e i valori ottimali per ogni metrica e il processo a cui fanno riferimento.
\subsubsection{Elenco delle Metriche di Qualità} 
L'elenco delle \glossario{metriche} di qualità da adottare per il progetto è dettagliato secondo la struttura definita e copre diverse aree rilevanti per il processo di accertamento della qualità.
\paragraph{Metriche di processo}
\textbf{Planned Value}
\begin{itemize}
    \item \textbf{Notazione}: M.PC.PV
    \item \textbf{Descrizione}: rappresenta il costo stimato del lavoro programmato entro un determinato momento del progetto, secondo il piano di progetto originale
    \item \textbf{Formula}: 
    \begin{equation}
        BAC \times (\% Lavoro Pianificato)
    \end{equation}
    Dove:
    \begin{itemize}
        \item $BAC$ (\glossario{Budget at Completion}): budget previsto per la realizzazione del progetto
    \end{itemize}
    
\end{itemize}

\textbf{Earned Value}
\begin{itemize}
    \item \textbf{Notazione}: M.PC.EV
    \item \textbf{Descrizione}: rappresenta il valore del lavoro effettivamente completato alla data corrente
    \item \textbf{Formula}: 
    \begin{equation}
        BAC \times (\% Lavoro Completato)
    \end{equation}
    Dove:
    \begin{itemize}
        \item $BAC$(\glossario{Budget at Completion}): budget previsto per la realizzazione del progetto
    \end{itemize}
        
    
\end{itemize}
\textbf{Actual Cost}
\begin{itemize}
    \item \textbf{Notazione}: M.PC.AC
    \item \textbf{Descrizione}: rappresenta il costo sostenuto fino alla data corrente
    \item \textbf{Formula}: costo speso per la realizzazione delle attività svolte fino data corrente
\end{itemize}

\textbf{Schedule Variance}
\begin{itemize}
    \item \textbf{Notazione}: M.PC.SV
    \item \textbf{Descrizione}: misura la variazione in percentuale tra il valore del lavoro effettivamente completato e il valore del lavoro pianificato.
    Indica quindi se il progetto è in ritardo o in anticipo rispetto a quanto preventivato
    \item \textbf{Formula}:
    \begin{equation}
        \frac{EV - PV}{EV} \times 100
    \end{equation}
\end{itemize}

\textbf{Cost Variance}
\begin{itemize}
    \item \textbf{Notazione}: M.PC.CV
    \item \textbf{Descrizione}: misura la variazione in percentuale tra il valore del lavoro completato e il costo effettivo sostenuto per tale lavoro.
    Indica quindi se il progetto sta rispettando il budget pianificato
    \item \textbf{Formula}: 
    \begin{equation}
        \frac{EV - AC}{EV} \times 100
    \end{equation}
\end{itemize}

\textbf{Variazione del piano}
\begin{itemize}
    \item \textbf{Notazione}: M.PC.VP
    \item \textbf{Descrizione}: misura la variazione in percentuale tra il numero di task pianificati per un determinato periodo di tempo e il numero di task realmente realizzati.
    Misura quindi la qualità della pianificazione del lavoro per un periodo di tempo.
    \item \textbf{Formula}: 
    \begin{equation}
         \frac{T_p - T_c}{T_p} \times 100
    \end{equation}
    Dove:
    \begin{itemize}
        \item $T_p$: numero di task pianificati
        \item $T_c$: numero di task completati
    \end{itemize}
\end{itemize}

\textbf{Estimated at Completion}
\begin{itemize}
    \item \textbf{Notazione}: M.PC.EAC
    \item \textbf{Descrizione}: rappresenta una stima aggiornata del costo totale previsto per la realizzazione del progetto
    \item \textbf{Formula}: 
    \begin{equation}
        AC + (BAC - EV)
    \end{equation}
    Dove:
    \begin{itemize}
        \item $BAC$(\glossario{Budget at Completion}): budget previsto per la realizzazione del progetto.
    \end{itemize}
\end{itemize}

\textbf{Rischi inattesi}
\begin{itemize}
    \item \textbf{Notazione}: M.PC.RI
    \item \textbf{Descrizione}: indica il numero dei rischi che si sono verificati in un determinato periodo e che non erano stati preventivati tramite l'Analisi dei Rischi
    \item \textbf{Formula}: Numero di rischi occorsi e non preventivati
\end{itemize}
\textbf{Risk Mitigation Rate}
\begin{itemize}
    \item \textbf{Notazione}: M.PC.RMR
    \item \textbf{Descrizione}: indica la percentuale dei rischi occorsi durante lo sviluppo in un determinato periodo che sono stati mitigati correttamente tramite le strategie di mitigazione indicate nella Analisi dei Rischi
    \item \textbf{Formula}: 
    \begin{equation}
        \frac{R_g}{R_t} \times 100
    \end{equation}
    Dove:
    \begin{itemize}
        \item $R_g$: numero di rischi gestiti correttamente
        \item $R_t$: numero totale di rischi verificati nel periodo in analisi
    \end{itemize}
\end{itemize}

\textbf{Metriche Soddisfatte}
\begin{itemize}
    \item \textbf{Notazione}: M.PC.MS
    \item \textbf{Descrizione}: indica la percentuale di \glossario{metriche} soddisfatte, cioè con un valore calcolato che rispetta il valore minimo ammissibile indicato all'interno del Piano di Qualifica
    \item \textbf{Formula}:
    \begin{equation}
        \frac{M_s}{M_t} \times 100
    \end{equation}
    Dove:
    \begin{itemize}
        \item $M_s$: numero di \glossario{metriche} soddisfatte
        \item $M_t$: numero totale delle \glossario{metriche} analizzate
    \end{itemize}
\end{itemize}

\paragraph{Metriche di prodotto}
\textbf{Percentuale Requisiti Obbligatori Soddisfatti}
\begin{itemize}
    \item \textbf{Notazione}: M.PR.PRM
    \item \textbf{Descrizione}: rappresenta la percentuale di requisiti obbligatori soddisfatti rispetto i requisiti obbligatori totali inseriti all'interno del documento Analisi dei Requisiti
    \item \textbf{Caratteristiche}: Funzionalità
    \item \textbf{Formula}: 
    \begin{equation}
        \frac{RMS}{RMT} \times 100
    \end{equation}
    Dove:
    \begin{itemize}
        \item $RMS$: numero di requisiti obbligatori soddisfatti
        \item $RMT$: numero di requisiti obbligatori totale
    \end{itemize} 
\end{itemize}

\textbf{Percentuale Requisiti Opzionali Soddisfatti}
\begin{itemize}
    \item \textbf{Notazione}: M.PR.PRO
    \item \textbf{Descrizione}: rappresenta la percentuale di requisiti opzionali soddisfatti rispetto al numero totale di requisiti opzionali inseriti all'interno del documento di Analisi dei Requisiti
    \item \textbf{Caratteristiche}: Funzionalità
    \item \textbf{Formula}:  
    \begin{equation}
        \frac{ROS}{ROT} \times 100
    \end{equation}
    Dove:
    \begin{itemize}
        \item $ROS$: numero di requisiti opzionali soddisfatti
        \item $ROT$: numero totale di requisiti opzionali  
    \end{itemize}
\end{itemize}

\textbf{Correttezza Ortografica}
\begin{itemize}
    \item \textbf{Notazione}: M.PR.CO
    \item \textbf{Descrizione}: rappresenta il numero di errori ortografici rilevato all'interno di un documento 
    \item \textbf{Caratteristiche}: Usabilità
    \item \textbf{Formula}: N° errori ortografici rilevati
\end{itemize}

\textbf{Code Coverage}
\begin{itemize}
    \item \textbf{Notazione}: M.PR.CC
    \item \textbf{Descrizione}: misura la percentuale di codice sorgente che è stata eseguita durante i test;
    \item \textbf{Caratteristiche}: Manutenibilità
    \item \textbf{Formula}: 
    \begin{equation}
        \frac{L_{test}}{L_{tot}} \times 100
    \end{equation}
    Dove:
    \begin{itemize}
        \item $L_{test}$: numero di linee di codice testate
        \item $L_{tot}$: numero totale di linee di codice  
    \end{itemize}
\end{itemize}

\textbf{Branch Coverage}
\begin{itemize}
    \item \textbf{Notazione}: M.PR.BC
    \item \textbf{Descrizione}: Misura la percentuale di rami del codice (istruzioni di tipo condizionale come \textit{if, else, switch, ecc.}) che sono stati testati.
    \item \textbf{Caratteristiche}: Manutenibilità
    \item \textbf{Formula}: 
    \begin{equation}
        \frac{R_{test}}{R_{tot}} \times 100
    \end{equation}
    Dove:
    \begin{itemize}
        \item $R_{test}$: numero di rami eseguiti durante i test
        \item $R_{tot}$: numero totale di rami nel codice 
    \end{itemize}
\end{itemize}

\textbf{Complessità ciclomatica}
\begin{itemize}
    \item \textbf{Notazione}: M.PR.COC
    \item \textbf{Descrizione}: Indica la complessità del codice, calcolata sulla base del numero di percorsi indipendenti in una esecuzione con singolo ingresso e singola uscita
    \item \textbf{Caratteristiche}: Manutenibilità
    \item \textbf{Formula}: 
    \begin{equation}
        E - N + 2P
    \end{equation}
    Dove:
    \begin{itemize}
        \item $E$: numero di archi nel grafico di flusso di controllo
        \item $N$: numero di nodi nel grafico di flusso di controllo
        \item $P$: numero delle componenti connesse da ogni arco
    \end{itemize}
\end{itemize}

\textbf{Tempo di risposta}
\begin{itemize}  
    \item \textbf{Notazione}: M.PR.TR
    \item \textbf{Descrizione}:  Il tempo impiegato dal sistema per rispondere a una richiesta.
    \item \textbf{Caratteristiche}: Efficienza
    \item \textbf{Formula}: tempo che intercorre tra la richiesta dell'esecuzione di un test su un dataset e l'ottenimento dei risultati
\end{itemize}

\textbf{Browser supportati}
\begin{itemize}  
    \item \textbf{Notazione}: M.PR.BS
    \item \textbf{Descrizione}:  percentuale di browser su cui l'applicazione risulta utilizzabile
    \item \textbf{Caratteristiche}: Portabilità
    \item \textbf{Formula}: 
    \begin{equation}
        \frac{B_s}{B_t} \times 100
    \end{equation}
    Dove:
    \begin{itemize}
        \item $B_s$: numero di Browser supportati dall'applicazione
        \item $B_t$: numero totale di Browser
    \end{itemize}
\end{itemize}

\textbf{Numero di linee medie di codice per metodo}
\begin{itemize}  
    \item \textbf{Notazione}: M.PR.LCPM
    \item \textbf{Descrizione}:  lunghezza media in termine di linee di codice per metodi o funzioni di una classe
    \item \textbf{Caratteristiche}: Manutenibilità
    \item \textbf{Formula}: calcolato tramite uno script in \textit{Python}
\end{itemize}

\textbf{Percentuale di test superati}
\begin{itemize}  
    \item \textbf{Notazione}: M.PR.PTS
    \item \textbf{Descrizione}:  percentuale di test che hanno terminato con successo durante l'esecuzione
    \item \textbf{Caratteristiche}: Funzionalità
    \item \textbf{Formula}: 
    \begin{equation}
        \frac{T_s}{T_t} \times 100
    \end{equation}
    Dove:
    \begin{itemize}
        \item $T_s$: numero di test terminati con successo
        \item $T_t$: numero totale di test
    \end{itemize}
\end{itemize}

\textbf{Gestione degli errori}
\begin{itemize}  
    \item \textbf{Notazione}: M.PR.GE
    \item \textbf{Descrizione}:  percentuale di errori gestiti correttamente dall'applicazione
    \item \textbf{Caratteristiche}: Affidabilità
    \item \textbf{Formula}: 
    \begin{equation}
        \frac{E_g}{E_t} \times 100
    \end{equation}
    Dove:
    \begin{itemize}
        \item $E_g$: numero di errori gestiti correttamente 
        \item $E_t$: numero totale di errori previsti
    \end{itemize}
\end{itemize}

\subsubsection{Strumenti}
Qui vengono elencati gli strumenti usati, che vengono poi approfonditi nella sezione \hyperref[sec:strumenti]{Strumenti}.

\begin{itemize}
    \item \glossario{LaTeX};
    \item \glossario{Github};
    \item \glossario{Git};
\end{itemize}



